\subsection{主要结果的解读}
\label{sec:discussion-interpretation}

表~\ref{tab:denoising-metrics} 和表~\ref{tab:feature-preservation} 中的定量结果揭示了一个重要发现：PMDF-Net 与所有基线方法之间的性能差距不能归因于任何单一的架构或损失函数选择，而应归因于它们在统一多域框架内的有意结合。PMDF-Net 的 $+11.4$\,dB 信噪比(SNR)提升比最佳单域基线高出 $1.3$\,dB，这一绝对值差距看似不大，但代表着残余噪声方差降低超过 $35\%$。更重要的是，这一信噪比(SNR)增益并非以波形保真度为代价——PMDF-Net 同时达到了最低的到达时间误差 ($0.14$\,$\mu$s) 和最高的包络一致性 ($0.981$)，优于所有其他方法。这种去噪强度与特征保留的联合优化是该方法的核心特征。

实现这一双重成就的机制在于时域和时频分支所捕捉的互补信息。时域分支编码精确的时间动力学，包括到达时间和波包形态，而时频分支编码频率含量——当缺陷改变局部波速和厚度时，频率含量会发生特征性变化。单域架构必须将两个信息流压缩为一种表示；而多域融合则允许每个分支专门化。这种结构互补性解释了为何单域CNN尽管使用相同的物理约束损失函数进行训练，仍无法匹配 PMDF-Net 的性能——它缺乏分别推理时间起始和频谱偏差的建筑学词汇。

物理约束损失不仅仅是一个正则化器，更是一种必要的归纳偏置，防止网络将信噪比(SNR)指标作为信号失真的代理。若无特征保留项，仅使用均方误差(MSE)的目标函数会驱使网络产生过度平滑的输出，在最小化像素级误差的同时，抹去临床医生和检测人员所依赖的诊断特征。消融研究（第~\ref{sec:results-ablation} 节）证实了这一点：移除物理约束会导致信噪比(SNR)和特征指标同时下降，包络相关系数(CCC)降至临床相关性阈值 $0.95$ 以下。因此，物理损失函数对于下游检测任务的实用性而言并非可选调参，而是必要条件。

从实际部署的角度来看，高信噪比(SNR)提升与最小特征失真的结合使得采集时间大幅缩短成为可能。使用传统信号平均达到相当的缺陷检测准确度需要 256 次平均；而 PMDF-Net 从 16 次平均中即可提供统计上等效的检测，代表着激光发射次数减少 $16\times$。采集时间的压缩直接转化为工业检测中的吞吐量提升，以及医疗应用中的患者曝光减少，使该方法不仅是一项信号处理进步，更是一项操作层面的改进。